% Generated draft: requires mathematical and source review.
\documentclass[11pt]{article}
\usepackage[margin=25mm]{geometry}
\usepackage{fontspec}
\setmainfont{Latin Modern Roman}
\usepackage{amsmath,amssymb,mathtools,unicode-math}
\setmathfont{Latin Modern Math}
\usepackage{enumitem,needspace,xurl,hyperref}
\setlength{\parindent}{0pt}
\setlength{\parskip}{5pt}
\setlength{\emergencystretch}{3em}
\newcommand{\N}{\mathbb{N}}
\newcommand{\Z}{\mathbb{Z}}
\newcommand{\Q}{\mathbb{Q}}
\newcommand{\R}{\mathbb{R}}
\newcommand{\C}{\mathbb{C}}
\newcommand{\F}{\mathbb{F}}
\newcommand{\A}{\mathbb{A}}
\newcommand{\PP}{\mathbb{P}}
\newcommand{\E}{\mathbb{E}}
\newcommand{\eps}{\varepsilon}
\newcommand{\dd}{\,\mathrm{d}}
\begin{document}
\section*{UnsolvedMath ID: 10900006}
\textbf{Source problem code:} catalogue entry
\subsection{Definitions and mathematical statement}
The eight Thurston model geometries in dimension three are
\[
 S^3,\quad \R^3,\quad \mathbb H^3,\quad S^2\times\R,\quad
 \mathbb H^2\times\R,\quad \mathrm{Nil},\quad \mathrm{Sol},\quad
 \widetilde{\mathrm{SL}_2(\R)}.
\]
Here $\mathrm{Nil}$ is the three-dimensional Heisenberg geometry, $\mathrm{Sol}$ is the standard solvable three-dimensional geometry, and the last model is the universal-cover geometry associated with $\mathrm{SL}_2(\R)$. Each is understood with its Thurston geometric structure, not merely as an underlying topological space.

The source calls a \emph{geometric transition} a continuous path of geometries that abruptly changes type at a limit. Schematically, write $\mathcal G\rightsquigarrow\mathcal H$ if a family of structures of type $\mathcal G$ for $t>0$ converges, in the chosen transition framework, to a structure of type $\mathcal H$ as $t\downarrow0$. This is a directed relation. The topology, ambient model, and permitted coordinate changes or rescalings are essential to its meaning and are not specified by the excerpt.

\textbf{Problem (Leitner).} Describe all geometric transitions between the eight Thurston geometries, including which directed limits $\mathcal G\rightsquigarrow\mathcal H$ occur and which cannot occur. More generally, develop criteria to determine when one geometry is a limit of another, and identify restrictions that any limiting geometry must satisfy. The request to understand all transitions may include more than the existence relation between model types; the equivalence convention for classifying distinct transition families is also unspecified.

The source stresses that constructing a transition proves existence, whereas proving that a limit cannot occur is generally harder. The registry summary reports projective-geometric results for limits of hyperbolic geometry: Euclidean, Nil, and Sol possibilities, and exclusions of $\mathbb H^2\times\R$ and $\widetilde{\mathrm{SL}_2(\R)}$. It also reports many hyperbolic/anti-de Sitter transitions, a related framework involving a geometry outside the eight listed above. These framework-dependent reports do not by themselves define the unrestricted limiting relation in the question.
\subsection{Short English statement}
Classify the possible transitions among Thurston's eight three-dimensional geometries and find general tests and obstructions for one geometry to arise as a limit of another. The notion of convergence and allowed rescaling must be fixed.
\subsection{Sources}
\begin{itemize}
\item Local UnsolvedMath registry, permanent ID 10900006, source code catalogue entry; attributed to Leitner, with no proposal year supplied.
\item Delp, Hoffoss and Manning, \emph{Problems In Groups, Geometry, and Three-Manifolds} (2015), subsection 2.1 (Leitner), PDF page 2; supplied locator \url{https://arxiv.org/pdf/1512.04620#page=2}, and \url{https://arxiv.org/abs/1512.04620}. The provenance records source-TeX extraction checked against rendered PDF text and access on 2026-07-29.
\item The registry research summary reports projective-geometric limit classifications and hyperbolic/anti-de Sitter transitions without supplying individual primary references or full hypotheses.
\end{itemize}
\end{document}
